\documentclass[10pt,a4paper]{amsart}
\usepackage{amsmath,amssymb,amsthm,graphicx}
\usepackage{verbatim}
\newenvironment{reference}{\comment}{\endcomment}
%\newenvironment{reference}{}{}
\newenvironment{slogan}{\comment}{\endcomment}
\newenvironment{history}{\comment}{\endcomment}

% For commutative diagrams we use Xy-pic
\usepackage[all]{xy}

% We use 2cell for 2-commutative diagrams.
\xyoption{2cell}
% Theorem environments.
%
\theoremstyle{plain}
\newtheorem{theorem}[subsection]{Theorem}
\newtheorem{proposition}[subsection]{Proposition}
\newtheorem{lemma}[subsection]{Lemma}

\theoremstyle{definition}
\newtheorem{definition}[subsection]{Definition}
\newtheorem{example}[subsection]{Example}
\newtheorem{exercise}[subsection]{Exercise}
\newtheorem{situation}[subsection]{Situation}

\theoremstyle{remark}
\newtheorem{remark}[subsection]{Remark}
\newtheorem{remarks}[subsection]{Remarks}

\numberwithin{equation}{subsection}

% Macros
%
\def\lim{\mathop{\mathrm{lim}}\nolimits}
\def\colim{\mathop{\mathrm{colim}}\nolimits}
\def\Spec{\mathop{\mathrm{Spec}}}
\def\Hom{\mathop{\mathrm{Hom}}\nolimits}
\def\Ext{\mathop{\mathrm{Ext}}\nolimits}
\def\SheafHom{\mathop{\mathcal{H}\!\mathit{om}}\nolimits}
\def\SheafExt{\mathop{\mathcal{E}\!\mathit{xt}}\nolimits}
\def\Sch{\mathit{Sch}}
\def\Mor{\mathop{\mathrm{Mor}}\nolimits}
\def\Ob{\mathop{\mathrm{Ob}}\nolimits}
\def\Sh{\mathop{\mathit{Sh}}\nolimits}
\def\NL{\mathop{N\!L}\nolimits}
\def\CH{\mathop{\mathrm{CH}}\nolimits}
\def\proetale{{pro\text{-}\acute{e}tale}}
\def\etale{{\acute{e}tale}}
\def\QCoh{\mathit{QCoh}}
\def\Ker{\mathop{\mathrm{Ker}}}
\def\Im{\mathop{\mathrm{Im}}}
\def\Coker{\mathop{\mathrm{Coker}}}
\def\Coim{\mathop{\mathrm{Coim}}}

% Boxtimes
%
\DeclareMathSymbol{\boxtimes}{\mathbin}{AMSa}{"02}

%
% Macros for moduli stacks/spaces
%
\def\QCohstack{\mathcal{QC}\!\mathit{oh}}
\def\Cohstack{\mathcal{C}\!\mathit{oh}}
\def\Spacesstack{\mathcal{S}\!\mathit{paces}}
\def\Quotfunctor{\mathrm{Quot}}
\def\Hilbfunctor{\mathrm{Hilb}}
\def\Curvesstack{\mathcal{C}\!\mathit{urves}}
\def\Polarizedstack{\mathcal{P}\!\mathit{olarized}}
\def\Complexesstack{\mathcal{C}\!\mathit{omplexes}}
% \Pic is the operator that assigns to X its picard group, usage \Pic(X)
% \Picardstack_{X/B} denotes the Picard stack of X over B
% \Picardfunctor_{X/B} denotes the Picard functor of X over B
\def\Pic{\mathop{\mathrm{Pic}}\nolimits}
\def\Picardstack{\mathcal{P}\!\mathit{ic}}
\def\Picardfunctor{\mathrm{Pic}}
\def\Deformationcategory{\mathcal{D}\!\mathit{ef}}
\usepackage{cleveref}
\crefname{originalequation}{original equation}{original equations}
\Crefname{originalequation}{Original equation}{Original equations}
\makeatletter
\newlabel{equation-epsilon}{{08H7}{source}{}{sourcepointer-1}{}}
\newlabel{equation-epsilon@cref}{{[originalequation][0][]08H7}{[1][0][]source}{}{}{}}
\makeatother
\crefname{originalreference}{original reference}{original references}
\Crefname{originalreference}{Original reference}{Original references}
\makeatletter
\newlabel{lemma-derived-quasi-coherent-small-etale-site}{{071Q}{source}{}{sourcepointer-2}{}}
\newlabel{lemma-derived-quasi-coherent-small-etale-site@cref}{{[originalreference][0][]071Q}{[1][0][]source}{}{}{}}
\makeatother
\crefname{originalreference}{original reference}{original references}
\Crefname{originalreference}{Original reference}{Original references}
\makeatletter
\newlabel{lemma-descend-finite-type}{{08HD}{source}{}{sourcepointer-3}{}}
\newlabel{lemma-descend-finite-type@cref}{{[originalreference][0][]08HD}{[1][0][]source}{}{}{}}
\makeatother
\crefname{originalreference}{original reference}{original references}
\Crefname{originalreference}{Original reference}{Original references}
\makeatletter
\newlabel{lemma-descend-pseudo-coherent}{{08HE}{source}{}{sourcepointer-4}{}}
\newlabel{lemma-descend-pseudo-coherent@cref}{{[originalreference][0][]08HE}{[1][0][]source}{}{}{}}
\makeatother
\crefname{originalreference}{original reference}{original references}
\Crefname{originalreference}{Original reference}{Original references}
\makeatletter
\newlabel{lemma-epsilon-flat}{{08H8}{source}{}{sourcepointer-5}{}}
\newlabel{lemma-epsilon-flat@cref}{{[originalreference][0][]08H8}{[1][0][]source}{}{}{}}
\makeatother
\crefname{originallemma}{original lemma}{original lemmas}
\Crefname{originallemma}{Original lemma}{Original lemmas}
\makeatletter
\newlabel{perfect-lemma-identify-pseudo-coherent-noetherian}{{08E8}{source}{}{sourcepointer-6}{}}
\newlabel{perfect-lemma-identify-pseudo-coherent-noetherian@cref}{{[originallemma][0][]08E8}{[1][0][]source}{}{}{}}
\makeatother
\crefname{originallemma}{original lemma}{original lemmas}
\Crefname{originallemma}{Original lemma}{Original lemmas}
\makeatletter
\newlabel{sites-cohomology-lemma-pseudo-coherent-independent-representative}{{08FU}{source}{}{sourcepointer-7}{}}
\newlabel{sites-cohomology-lemma-pseudo-coherent-independent-representative@cref}{{[originallemma][0][]08FU}{[1][0][]source}{}{}{}}
\makeatother
\crefname{originallemma}{original lemma}{original lemmas}
\Crefname{originallemma}{Original lemma}{Original lemmas}
\makeatletter
\newlabel{spaces-cohomology-lemma-coherent-Noetherian}{{07UB}{source}{}{sourcepointer-8}{}}
\newlabel{spaces-cohomology-lemma-coherent-Noetherian@cref}{{[originallemma][0][]07UB}{[1][0][]source}{}{}{}}
\makeatother
\crefname{originallemma}{original lemma}{original lemmas}
\Crefname{originallemma}{Original lemma}{Original lemmas}
\makeatletter
\newlabel{spaces-properties-lemma-quasi-compact-affine-cover}{{03H6}{source}{}{sourcepointer-9}{}}
\newlabel{spaces-properties-lemma-quasi-compact-affine-cover@cref}{{[originallemma][0][]03H6}{[1][0][]source}{}{}{}}
\makeatother
\title{On a Noetherian space, m-pseudo-coherence is coherence of higher cohomology sheaves}
\author{The Stacks Project Authors}
\date{stacks-project:spaces-perfect.tex:local-snapshot-20261012; GNU FDL 1.2 or later}
\begin{document}
\maketitle
\noindent\textit{Editorial note.} All delivered mathematical text is the author's own, copied from the pinned Stacks Project chapter file spaces-perfect.tex; no statement or proof step was rewritten, reordered or removed.\par\medskip
\noindent\textit{Editorial note.} The counted proof is the author's complete argument: the proof reduces to the representable case by etale descent of both sides, then transfers m-pseudo-coherence, coherence and vanishing along epsilon\textasciicircum{}* and concludes by the scheme-case lemma of Derived Categories of Schemes\par\medskip
\noindent\textit{Editorial note.} The label lemma-identify-pseudo-coherent-noetherian names the counted statement, and the proof environment immediately following it in the source is delivered verbatim as its proof.\par\medskip
\noindent\textit{Editorial note.} The counted statement is carried by the author's own lemma environment in the chapter source.\par\medskip
\noindent\textit{Allowed standard prerequisites (editorial).}\begin{itemize}
\item quasi-compact algebraic spaces admit surjective etale maps from affine schemes (Properties of Spaces, Lemma spaces-properties-lemma-quasi-compact-affine-cover)
\item m-pseudo-coherence is insensitive to replacing the space by an etale cover (Cohomology on Sites, Lemma sites-cohomology-lemma-pseudo-coherent-independent-representative)
\item coherence of sheaves is etale local on Noetherian spaces (Cohomology of Spaces, Lemma spaces-cohomology-lemma-coherent-Noetherian)
\item the comparison of D\_QCoh on the small etale site with the Zariski site of a scheme via epsilon, with E = epsilon\textasciicircum{}*E\_0 (lemma-derived-quasi-coherent-small-etale-site of the same chapter)
\item descent of m-pseudo-coherence, of finite type, and the flatness of epsilon (lemma-descend-pseudo-coherent, lemma-descend-finite-type and lemma-epsilon-flat of the same chapter)
\item the scheme case: on a Noetherian scheme, m-pseudo-coherence is equivalent to coherence of H\textasciicircum{}i for i >= m and vanishing for i >> 0 (Derived Categories of Schemes, Lemma perfect-lemma-identify-pseudo-coherent-noetherian)
\end{itemize}\medskip
\ifdefined\hypertarget\hypertarget{sourcepointer-1}{}\fi
\noindent\textit{Source reference.} Original equation 08H7: sources/stacks-spaces-perfect/pinned-original-source.tex, source lines 320--320. Prerequisite pointer invoked at proof line(s) 3356; it names a standard fact of the same chapter that the counted argument uses, and every step of the counted proof built on it is retained verbatim in the delivered body.\par\medskip
\ifdefined\hypertarget\hypertarget{sourcepointer-2}{}\fi
\noindent\textit{Source reference.} Original reference 071Q: sources/stacks-spaces-perfect/pinned-original-source.tex, source lines 402--402. Prerequisite pointer invoked at proof line(s) 3352; it names a standard fact of the same chapter that the counted argument uses, and every step of the counted proof built on it is retained verbatim in the delivered body.\par\medskip
\ifdefined\hypertarget\hypertarget{sourcepointer-3}{}\fi
\noindent\textit{Source reference.} Original reference 08HD: sources/stacks-spaces-perfect/pinned-original-source.tex, source lines 3081--3081. Prerequisite pointer invoked at proof line(s) 3360; it names a standard fact of the same chapter that the counted argument uses, and every step of the counted proof built on it is retained verbatim in the delivered body.\par\medskip
\ifdefined\hypertarget\hypertarget{sourcepointer-4}{}\fi
\noindent\textit{Source reference.} Original reference 08HE: sources/stacks-spaces-perfect/pinned-original-source.tex, source lines 3133--3133. Prerequisite pointer invoked at proof line(s) 3358; it names a standard fact of the same chapter that the counted argument uses, and every step of the counted proof built on it is retained verbatim in the delivered body.\par\medskip
\ifdefined\hypertarget\hypertarget{sourcepointer-5}{}\fi
\noindent\textit{Source reference.} Original reference 08H8: sources/stacks-spaces-perfect/pinned-original-source.tex, source lines 330--330. Prerequisite pointer invoked at proof line(s) 3362; it names a standard fact of the same chapter that the counted argument uses, and every step of the counted proof built on it is retained verbatim in the delivered body.\par\medskip
\ifdefined\hypertarget\hypertarget{sourcepointer-6}{}\fi
\noindent\textit{Source reference.} Original lemma 08E8: sources/stacks-chow/tags.txt, source lines 10664--10664. Prerequisite pointer invoked at proof line(s) 3365; it names a standard fact of a sibling chapter that the counted argument uses, and every step of the counted proof built on it is retained verbatim in the delivered body.\par\medskip
\ifdefined\hypertarget\hypertarget{sourcepointer-7}{}\fi
\noindent\textit{Source reference.} Original lemma 08FU: sources/stacks-chow/tags.txt, source lines 10722--10722. Prerequisite pointer invoked at proof line(s) 3344; it names a standard fact of a sibling chapter that the counted argument uses, and every step of the counted proof built on it is retained verbatim in the delivered body.\par\medskip
\ifdefined\hypertarget\hypertarget{sourcepointer-8}{}\fi
\noindent\textit{Source reference.} Original lemma 07UB: sources/stacks-chow/tags.txt, source lines 9966--9966. Prerequisite pointer invoked at proof line(s) 3347; it names a standard fact of a sibling chapter that the counted argument uses, and every step of the counted proof built on it is retained verbatim in the delivered body.\par\medskip
\ifdefined\hypertarget\hypertarget{sourcepointer-9}{}\fi
\noindent\textit{Source reference.} Original lemma 03H6: sources/stacks-chow/tags.txt, source lines 4555--4555. Prerequisite pointer invoked at proof line(s) 3339; it names a standard fact of a sibling chapter that the counted argument uses, and every step of the counted proof built on it is retained verbatim in the delivered body.\par\medskip
\begin{lemma}
\label{lemma-identify-pseudo-coherent-noetherian}
Let $S$ be a scheme. Let $X$ be a Noetherian algebraic space over $S$.
Let $E$ be an object of $D_\QCoh(\mathcal{O}_X)$. For
$m \in \mathbf{Z}$ the following are equivalent
\begin{enumerate}
\item $H^i(E)$ is coherent for $i \geq m$ and zero for $i \gg 0$, and
\item $E$ is $m$-pseudo-coherent.
\end{enumerate}
In particular, $E$ is pseudo-coherent if and only if $E$ is an object
of $D^-_{\textit{Coh}}(\mathcal{O}_X)$.
\end{lemma}

\begin{proof}
As $X$ is quasi-compact we can find an affine scheme $U$ and a surjective
\'etale morphism $U \to X$ (Properties of Spaces, Lemma
\ref{spaces-properties-lemma-quasi-compact-affine-cover}).
Observe that $U$ is Noetherian.
Note that $E$ is $m$-pseudo-coherent if and only if $E|_U$ is
$m$-pseudo-coherent (follows from the definition or from
Cohomology on Sites, Lemma
\ref{sites-cohomology-lemma-pseudo-coherent-independent-representative}).
Similarly, $H^i(E)$ is coherent if and only if $H^i(E)|_U = H^i(E|_U)$
is coherent (see Cohomology of Spaces, Lemma
\ref{spaces-cohomology-lemma-coherent-Noetherian}).
Thus we may assume that $X$ is representable.

\medskip\noindent
If $X$ is representable by a scheme $X_0$ then
(Lemma \ref{lemma-derived-quasi-coherent-small-etale-site})
we can write $E = \epsilon^*E_0$ where $E_0$ is an object of
$D_\QCoh(\mathcal{O}_{X_0})$ and
$\epsilon : X_\etale \to (X_0)_{Zar}$ is as in
(\ref{equation-epsilon}).
In this case $E$ is $m$-pseudo-coherent
if and only if $E_0$ is by Lemma \ref{lemma-descend-pseudo-coherent}.
Similarly, $H^i(E_0)$ is of finite type (i.e., coherent) if and only if
$H^i(E)$ is by Lemma \ref{lemma-descend-finite-type}.
Finally, $H^i(E_0) = 0$ if and only if $H^i(E) = 0$ by
Lemma \ref{lemma-epsilon-flat}.
Thus we reduce to the case of schemes which is
Derived Categories of Schemes, Lemma
\ref{perfect-lemma-identify-pseudo-coherent-noetherian}.
\end{proof}

\end{document}
